\problemname{Tina Scoreboard}
ICPC is a programming competition where you have five hours to solve as many problems as possible.
This problem is about a competition that is similar, but not exactly the same as the official ICPC.

In this competition, teams make submissions on problems.
For each submission, you do not receive points: you only get Accepted or Wrong Answer.

At the end of the contest, a ranking of the teams is established in the following way.
First, we count how many problems each team has solved; let this be denoted by an integer $p$.
Then we compute each team's \textit{submission penalty}.
For each problem that a team gets Accepted on, the submission penalty is equal to the number of
incorrect submissions the team made on that problem.
The total submission penalty for a team is the sum of the submission penalties over all problems
\underline{that they solved}, which we denote by an integer $t$.
Each team's result can therefore be described by the ordered pair $(p, t)$.
Note that if a team never solves a certain problem during the entire contest, then they receive
no submission penalty from that problem.

The rank of a given team is the number of teams that are \textit{strictly better}.
A team $i$ is strictly better than another team $j$ if either
\[
(p_i > p_j) \text{ or } (p_i = p_j \text{ and } t_i < t_j)
\]
holds.

The table below shows an example of the ranks that different teams receive, given a set of $p$
and $t$ values.

\begin{center}
\begin{tabular}{| c | c | c |}
  \hline
  Rank & $p$ & $t$ \\ \hline
  $0$ & $10$ & $500$ \\ \hline
  $1$ & $9$ & $1000$ \\ \hline
  $2$ & $8$ & $500$ \\ \hline
  $3$ & $7$ & $300$ \\ \hline
  $3$ & $7$ & $300$ \\ \hline
  $5$ & $6$ & $250$ \\ \hline
\end{tabular}
\end{center}

Note that if two teams have exactly the same result, they receive the same rank.

During the first four hours of the contest, each team can see whether their submissions were Accepted.
They can also see the results of other teams' submissions via the scoreboard.

During the last hour, the scoreboard is \textit{frozen}, and thus teams can only see the results
of their own submissions.
Each team can still see when new submissions are made by other teams, but not their results.

When the contest is over, it is time to \textit{unfreeze} the scoreboard.
This is done by revealing the result of each submission one at a time.
Depending on the order in which this is done, it can be more or less exciting for the teams.
This is because the rank of a team is primarily determined by how many problems they have solved.
For example, if a certain team has solved $4$ problems and it is revealed that all other teams
have solved more, then that team loses interest.
If the submissions had instead been revealed in an order such that it takes longer for them to
find this out, it would of course be more exciting.

Tina aims to make ICPC more exciting.
She has therefore proposed an order in which the results of submissions should be revealed.
Given this order, she wants to compute for each team the number of submissions that are revealed
before they lose interest.

More precisely, assume that after each revealed submission, every team considers all possible
outcomes of finishing the unfreezing of the scoreboard.
If a given team has the same final position in all these outcomes, then that team loses interest.
Tina then wants to know, for each team, the number of submissions that were revealed before they
lost interest.

For the order to be reasonable, she guarantees that each team's submissions on a given problem
are ordered by time, so that earlier submissions come first.

\section*{Input}
The first line contains four integers $N$, $P$, $H$, and $F$
($1 \leq N \leq 2 \cdot 10^5$, $1 \leq P \leq 15$, $0 \leq H \leq 2 \cdot 10^5$, $1 \leq F \leq 2 \cdot 10^5$),
the number of teams, the number of problems, the number of submissions before the freeze,
and the number of submissions during the freeze.

The following $H$ lines describe the submissions that occurred before the freeze.
Each such submission is described by the integers $w$, $p$, and $r$
($1 \leq w \leq N$, $1 \leq p \leq P$, $r \in \{AC, WA\}$),
meaning that the team with index $w$ submitted to problem $p$, and the result was $r$.
If $r$ is AC, the submission was Accepted; otherwise it was not.
These submissions are given in chronological order.

The following $F$ lines describe the submissions that occurred during the freeze.
Their format is identical to the submissions that occurred before the freeze.
They are given in the order that Tina is curious about.
It is guaranteed that for each team, submissions to a given problem appear in chronological order.
No other guarantees are given.

\textbf{Note:} it is also guaranteed that a team will never submit to a problem that they have
already solved, either before or during the freeze.

\section*{Output}
For each team, output the number of submissions that were revealed before the team lost interest.
Print these values space-separated on a single line.
Note that if a team knows its final position before the unfreezing begins, you should output $0$
for that team.

\section*{Scoring}
Your solution will be tested on a set of test groups, each worth a number of points. Each test group contains
a set of test cases. To get the points for a test group you need to solve all test cases in the test group.

\begin{center}
\begin{tabular}{| l | l | p{12cm} |}
  \hline
  \textbf{Group} & \textbf{Points} & \textbf{Constraints} \\ \hline
  $1$ & $5$ & $N \leq 50$, $P = 1$, and each team makes at most one submission per problem. \\ \hline
  $2$ & $6$ & $N \leq 50$, and each team makes at most one submission per problem. \\ \hline
  $3$ & $15$ & $N, F, H \leq 50$ \\ \hline
  $4$ & $8$ & $N, F, H \leq 500$ \\ \hline
  $5$ & $30$ & $N, F, H \leq 2000$ \\ \hline
  $6$ & $36$ & No additional constraints. \\ \hline
\end{tabular}
\end{center}

\section*{Explanation of the Samples}
Sample case 1: before the scoreboard is frozen, team $2$ already knows that they will finish with
$1$ solved problem and $0$ submission penalty.
Thus, they already know their final rank.
Team $1$ does not know whether team $2$ will solve the problem or not, but even if the other team
solves the problem, their result cannot worsen the rank of team $1$.

Sample case 2: similarly to sample case $1$, team $2$ already has all the information.
Team $1$, however, cannot be certain of its final position immediately after the contest,
since if team $2$ solves the problem, team $1$ loses one rank.
When the result is finally revealed, they become certain.

Sample case 3: even though it may not seem so, team $1$ can immediately figure out that they have
won the contest.
Since a team never submits to a problem they have already solved, team $1$ knows that the only
submission from team $2$ that could be Accepted is the last one.
Even if that submission is Accepted, it will not affect their rank.
